\section*{Hoofdstuk \ref{ch:b2:plant-life-cycles} --- Levenscycli en voortplanting van landplanten}

\begin{solution}{exo:b2:plant-life-cycles:1}
\omterm{def:b2:plant-life-cycles:cycles}{Gametofyt}: de haploïde meercellige generatie, door mitose uit een spore
gegroeid, die door mitose gameten maakt. \omterm{def:b2:plant-life-cycles:cycles}{Sporofyt}: de diploïde
generatie, door mitose uit een zygote gegroeid, die door \omterm{def:b2:meiosis-heredity:meiosis}{meiose} sporen
maakt. Spore: een haploïde cel, door \omterm{def:b2:meiosis-heredity:meiosis}{meiose} gemaakt, die uitgroeit zonder
te versmelten. Gameet: een haploïde cel, (bij planten) door mitose
gemaakt, die met een andere moet versmelten.
\end{solution}

\begin{solution}{exo:b2:plant-life-cycles:2}
Bladcellen zijn $2n = 1260$: spore 630, cel van het \omterm{def:b2:plant-life-cycles:fern}{prothallium} 630,
zaadcel 630, eicel 630, zygote 1260.
\end{solution}

\begin{solution}{exo:b2:plant-life-cycles:3}
Groen kussen: \omterm{def:b2:plant-life-cycles:cycles}{gametofyt}. Steeltje met kapsel: \omterm{def:b2:plant-life-cycles:cycles}{sporofyt}. Varenblad:
\omterm{def:b2:plant-life-cycles:cycles}{sporofyt}. \omterm{def:b2:plant-life-cycles:fern}{Prothallium}: \omterm{def:b2:plant-life-cycles:cycles}{gametofyt}. Stuifmeelkorrel: mannelijke
\omterm{def:b2:plant-life-cycles:cycles}{gametofyt}. Voedselvoorraad van het zaad (bij een den): vrouwelijke
\omterm{def:b2:plant-life-cycles:cycles}{gametofyt}. Embryo van het zaad: de volgende \omterm{def:b2:plant-life-cycles:cycles}{sporofyt}.
\end{solution}

\begin{solution}{exo:b2:plant-life-cycles:4}
Haplontisch (\emph{Chlamydomonas}): \omterm{def:b2:meiosis-heredity:meiosis}{meiose} in de zygote, meteen.
Diplontisch (dieren, \emph{Fucus}): de \omterm{def:b2:meiosis-heredity:meiosis}{meiose} maakt de gameten.
Haplodiplontisch (mossen, \omterm{def:b2:plant-life-cycles:fern}{varens}, zaadplanten, veel algen): de \omterm{def:b2:meiosis-heredity:meiosis}{meiose}
maakt sporen in de sporangia van de \omterm{def:b2:plant-life-cycles:cycles}{sporofyt}.
\end{solution}

\begin{solution}{exo:b2:plant-life-cycles:5}
$0.03/10^{-4} = \qty{300}{s}$, vijf minuten. In \qty{20}{min}
\qty{12}{cm}. Bevruchting werkt alleen tussen planten die enkele
centimeters van elkaar staan, zodat mossen in dichte kussens groeien met
beide geslachten (of beide organen) binnen bereik.
\end{solution}

\begin{solution}{exo:b2:plant-life-cycles:6}
$v_s = 2\times(1.5\times 10^{-5})^2\times 1099\times 9.81/(9\times
1.8\times 10^{-5}) = \qty{0.030}{m/s}$, \qty{3}{cm/s}; afstand $uh/v_s
= 2\times 0.5/0.03 = \qty{33}{m}$. Zaad: $2\times 20/0.8 =
\qty{50}{m}$ --- een zwaardere eenheid, van hoger losgelaten en met een
vleugel, komt ongeveer even ver.
\end{solution}

\begin{solution}{exo:b2:plant-life-cycles:7}
$200\times 40\times 64 = \num{512000}$ sporen per blad; $5.1$
prothallia; $0.1$ jonge \omterm{def:b2:plant-life-cycles:cycles}{sporofyt} per blad.
\end{solution}

\begin{solution}{exo:b2:plant-life-cycles:8}
Een zaad is een behouden \omterm{def:b2:plant-life-cycles:heterospory}{zaadbeginsel}: de vrouwelijke \omterm{def:b2:plant-life-cycles:cycles}{gametofyt} moet op
de ouderplant blijven, ingesloten en gevoed, terwijl de mannelijke
\omterm{def:b2:plant-life-cycles:cycles}{gametofyt} ernaartoe moet reizen. Daarvoor zijn twee soorten sporen met
twee bestemmingen nodig. De enige spore van een isospore plant moet
worden uitgestrooid om uit te groeien tot een vrije tweeslachtige
\omterm{def:b2:plant-life-cycles:cycles}{gametofyt}; haar behouden zou ook de mannelijke functie behouden en elke
plant tot zelfbevruchting dwingen, en er zou niets te reizen overblijven.
\end{solution}

\begin{solution}{exo:b2:plant-life-cycles:9}
Spore: $\tfrac43\pi(1.5\times 10^{-5})^3\times 1100 =
\qty{1.6e-11}{kg}$, \qty{16}{ng}; energie $0.5\times 1.6\times
10^{-8}\times 2\times 10^{4} = \qty{1.6e-4}{J}$. Zaad: \qty{5}{mg},
$0.5\times 5\times 10^{-3}\times 2\times 10^{4} = \qty{50}{J}$ ---
$3\times 10^{5}$ keer meer. De spore kan één klein fotosynthetisch
schubje maken en niets vóór er licht is; het zaad kan in het donker een
wortel en een scheut maken, een seizoen wachten, en overleven als het
wordt opgegeten of uitgedroogd.
\end{solution}

\begin{solution}{exo:b2:plant-life-cycles:10}
$10^{11}/(10^{6}\times 20) = \num{5000}$ korrels per kubieke meter; flux
$5000\times 10^{-7}\times 2 = \qty{1e-3}{}$ korrels per seconde; ongeveer
\qty{1000}{s}, een kwartier, per korrel. De kegel moet uren tot dagen open
en ontvankelijk zijn, en precies wanneer de mannelijke kegels \omterm{def:b2:plant-life-cycles:heterospory}{stuifmeel}
afgeven --- de bestuiving is tot op de week getimed.
\end{solution}

\begin{solution}{exo:b2:plant-life-cycles:11}
Diploïdie maskeert recessieve schadelijke \omterm{def:b2:genome-diversification:mutation}{mutaties}; een \omterm{def:b2:meiosis-heredity:meiosis}{diploïd} lichaam
met vaatweefsel kan groot en langlevend zijn, en hoogte helpt zowel bij
het invangen van licht als bij de \omterm{thm:b2:plant-life-cycles:dispersal}{verspreiding van sporen}; sporen die hoog
op een \omterm{def:b2:plant-life-cycles:cycles}{sporofyt} worden gemaakt, verspreiden zich door de lucht, terwijl
gameten moeten zwemmen. De haploïde generatie blijft bestaan omdat \omterm{def:b2:meiosis-heredity:meiosis}{meiose}
en bevruchting een haploïde fase vereisen, en de stuifmeelkorrel en de
embryozak zijn de minimale lichamen die gameten naar elkaar brengen en het
embryo voeden.
\end{solution}

\begin{solution}{exo:b2:plant-life-cycles:12}
Door sporen te kweken en ze onder de microscoop te volgen, kon hij
vaststellen dat een spore uitgroeit tot een klein lichaam met
geslachtsorganen, dat het embryo in het \omterm{def:b2:plant-life-cycles:moss}{archegonium} uit de bevruchte eicel
ontstaat, dat de sporendragende plant daaruit groeit, en dat dezelfde
organen in gereduceerde vorm in het \omterm{def:b2:plant-life-cycles:heterospory}{zaadbeginsel} van een conifeer bestaan:
de \omterm{def:b2:plant-life-cycles:cycles}{generatiewisseling} als een opeenvolging van lichamen en organen. Hij kon
niet weten wat de twee generaties op het niveau van de cel onderscheidde.
Chromosoomtellingen toonden aan dat de twee lichamen een factor twee in
chromosoomaantal verschillen, dat de \omterm{def:b2:meiosis-heredity:meiosis}{meiose} bij de sporenvorming ligt en de
bevruchting bij de eicel, en verklaarden zo waarom de \omterm{def:b2:plant-life-cycles:cycles}{generatiewisseling}
verplicht is.
\end{solution}

\begin{solution}{pb:b2:plant-life-cycles:1}
\textbf{1.} $10\times 200\times 40\times 64 = 5.1\times 10^{6}$
sporen.
\textbf{2.} $\tfrac43\pi(1.5\times 10^{-5})^3\times 1100 =
\qty{1.6e-11}{kg}$ per spore; oogst $8\times 10^{-5}$ kg, \qty{80}{mg}.
\textbf{3.} $v_s = 2\times 2.25\times 10^{-10}\times 1099\times
9.81/(1.62\times 10^{-4}) = \qty{0.030}{m/s}$.
\textbf{4.} $x = 1.5\times 0.5/0.030 = \qty{25}{m}$; schijf van
$\pi\times 25^2 = \qty{2000}{m^2}$; ongeveer \num{2600} sporen per
vierkante meter.
\textbf{5.} Opwaartse luchtstromen van enkele centimeters per seconde
overtreffen $v_s$, zodat een spore die in turbulentie terechtkomt uren in
de lucht blijft en honderden kilometers aflegt; één enkele spore sticht
een populatie als haar \omterm{def:b2:plant-life-cycles:fern}{prothallium} zichzelf kan bevruchten (het draagt
beide organen), en daarom hebben afgelegen eilanden rijke varenflora's.
\textbf{6.} Spore: $0.5\times 1.6\times 10^{-8}\times 2\times 10^{4} =
\qty{1.6e-4}{J}$; oogst: \qty{820}{J}; één dennenzaad: \qty{50}{J} ---
de hele sporenoogst is zestien zaden waard.
\textbf{7.} $5.1\times 10^{6}/2\times 10^{4} = 256$ prothallia.
\textbf{8.} $1 - 0.8^{6} = 1 - 0.26 = 0.74$.
\textbf{9.} $256\times 0.74\times 0.5\times 0.02 = 1.9$ volwassen \omterm{def:b2:plant-life-cycles:fern}{varens}
per seizoen.
\textbf{10.} $30\times 1.9 \approx 57$: veel meer dan de ene vervanger van
een stabiele populatie, dus de overlevingscijfers zijn optimistisch ---
in een volle habitat sterven de meeste jonge \omterm{def:b2:plant-life-cycles:cycles}{sporofyten} door verdringing
en schaduw.
\textbf{11.} $10^{-3}/10^{-4} = \qty{10}{s}$. Veel \omterm{def:b2:plant-life-cycles:fern}{varens} laten
antheridia en archegonia op verschillende tijdstippen rijpen, of geven een
hormoon af (antheridiogeen) dat naburige jonge prothallia mannelijk maakt,
zodat de zaadcellen van een ander individu komen.
\textbf{12.} Het is één cel dik, zonder wortels, onbeschermd, heeft
vloeibaar water nodig voor de bevruchting en leeft enkele weken: vrijwel
alle sterfte van de cyclus valt erop (één spore op $2\times 10^{4}$ wordt
er een; een kwart daarvan maakt nooit een natte nacht mee).
\textbf{13.} $v_s = 2\times 6.25\times 10^{-10}\times 599\times
9.81/(1.62\times 10^{-4}) = \qty{0.045}{m/s}$.
\textbf{14.} $3\times 15/0.045 = \qty{1000}{m}$.
\textbf{15.} $10^{11}/(2\times 10^{7}) = \num{5000}$ per kubieke meter.
\textbf{16.} $5000\times 10^{-7}\times 2 = \qty{1e-3}{s^{-1}}$: 3.6 per
uur; de eerste komt na ongeveer \qty{1000}{s}, een kwartier.
\textbf{17.} Vijf kilometer benedenwinds heeft de wolk zich zijwaarts en
omhoog verspreid en zijn de meeste korrels neergedaald (een bereik van
\qty{1}{km} per \qty{15}{m} hoogte), zodat de concentratie ordes van
grootte lager is en een \omterm{def:b2:plant-life-cycles:heterospory}{zaadbeginsel} dagen op een korrel kan wachten; het
eigen \omterm{def:b2:plant-life-cycles:heterospory}{stuifmeel} van de boom geeft vooral zelfbestoven zaden, die aborteren.
Windbestuiving vraagt een dichte wolk, vandaar bestanden.
\textbf{18.} $10^{14}$ korrels voor $2\times 10^{6}$ \omterm{def:b2:plant-life-cycles:heterospory}{zaadbeginsels}: $5\times
10^{7}$ korrels per \omterm{def:b2:plant-life-cycles:heterospory}{zaadbeginsel}.
\textbf{19.} $2\times 10^{6}$ \omterm{def:b2:plant-life-cycles:heterospory}{zaadbeginsels} $\times 0.6\times 0.8 = 9.6\times
10^{5}$ zaden; \qty{4.8}{kg}; \qty{48}{MJ}.
\textbf{20.} $3\times 20/0.8 = \qty{75}{m}$: het \omterm{def:b2:plant-life-cycles:heterospory}{stuifmeel} komt een
kilometer ver, het zaad enkele boomhoogten; de genen reizen met het
\omterm{def:b2:plant-life-cycles:heterospory}{stuifmeel}, de plant met het zaad.
\textbf{21.} $9.6\times 10^{5}\times 0.05\times 0.01 = 480$ volwassen
nakomelingen, tegen 57 bij de \omterm{def:b2:plant-life-cycles:fern}{varen}.
\textbf{22.} Den: $4.8\times 10^{7}/480 = \qty{100}{kJ}$ per volwassen
nakomeling; \omterm{def:b2:plant-life-cycles:fern}{varen}: $820/1.9 = \qty{430}{J}$ per volwassen nakomeling,
tweehonderd keer minder.
\textbf{23.} Beide bestaan voor de helft uit droge stof à \qty{20}{kJ/g},
zodat bij een basaal verbruik van, zeg, \qty{0.5}{mW} per gram droge stof
beide $10^{4}\,
\text{J/g}/(5\times 10^{-4}\,\text{W/g}) = 2\times 10^{7}$ s meegaan,
ongeveer acht maanden: de duur van de rustperiode is niet wat de energie
van het zaad koopt. Ze koopt wat er na de kieming gebeurt --- een wortel
en een scheut, in het donker uit reserves opgebouwd, weken voordat de
kiemplant zichzelf voedt --- en, met de zaadhuid, bescherming in de
tussentijd; de \qty{1.6e-4}{J} van de spore bouwt enkele cellen die
meteen moeten fotosynthetiseren.
\textbf{24.} Zaad: bevruchting zonder water; een bevoorraad, beschermd,
rustend embryo dat zich op droge of seizoensgebonden plaatsen vestigt;
verspreiding door vleugels, vruchten en dieren. De strategie van de \omterm{def:b2:plant-life-cycles:fern}{varen}
wint in natte, schaduwrijke, stabiele habitats, en bij de kolonisatie van
verre of pas vrijgekomen grond, waar goedkope, lichte, talrijke
verspreidingseenheden meer tellen dan bevoorrading.
\textbf{25.} \omterm{def:b2:plant-life-cycles:fern}{Varen}: $5.1\times 10^{6}$ sporen per seizoen voor 1.9
volwassenen, $2.7\times 10^{6}$ sporen per volwassen nakomeling,
\qty{430}{J} per stuk; den: $9.6\times 10^{5}$ zaden voor 480 volwassenen,
\num{2000} zaden per volwassen nakomeling, \qty{100}{kJ} per stuk.
\end{solution}
